\section{Literature review}

Our proposed research focuses on the internal geometry of LLMs and the mechanics of AI interpretability. Because the theoretical definitions of our problem space have been established, we ground our methodology directly in the following empirical works.

\paragraph{Universal and Transferable Adversarial Attacks on Aligned Language Models (Zou et al., 2023).}
Current methods for aligning LLMs to be safe rely heavily on behavioral fine-tuning, such as Reinforcement Learning from Human Feedback. However, \citet{zou2023gcg} demonstrated that these surface-level guardrails are fundamentally brittle. They introduced the Greedy Coordinate Gradient search method, which automatically generates adversarial prompt suffixes. When appended to a malicious query, these suffixes universally jailbreak aligned models by maximizing the probability of the model producing an affirmative response, entirely bypassing the safety training. This brittleness strongly motivates our project. Because behavioral filters can be bypassed so easily via input perturbations, we hypothesize that true safety alignment must be understood and engineered at the mechanistic level within the internal representations of the model itself.

\paragraph{The Linear Representation Hypothesis and the Geometry of Large Language Models (Park et al., 2023).}
While the theoretical definition of the Linear Representation Hypothesis (LRH) dictates that concepts occupy linear subspaces, \citet{park2023LRH} provided the empirical blueprint for testing and manipulating these spaces. They demonstrated that causal interventions on these vectors predictably alter generated text. By utilizing counterfactual input pairs to construct probes, they proved that geometric concepts such as cosine similarity and orthogonal projection map directly to causal changes in model behavior. Our project relies heavily on their methodology for causal separability. We apply their geometric framework to test whether intervening on an isolated safety vector mathematically leaves the utility capabilities of the model untouched, proving orthogonal separation.

\paragraph{Refusal in Language Models Is Mediated by a Single Direction (Arditi et al., 2024).}
Investigating how safety aligns with internal geometry, \citet{arditi2024refusal} discovered that an LLM's refusal behavior against malicious prompts is highly localized. It is not a distributed behavior scattered across the network but is instead mediated by a single, one-dimensional subspace. They isolated this vector using a Difference of Means (DoM) approach. By capturing the residual stream activations across layers and subtracting the mean activations of harmless instructions from those of harmful instructions, they found a ``refusal direction.'' Crucially, mathematically erasing this direction from the residual stream prevented the model from refusing harmful instructions entirely. We directly replicate this DoM methodology as our primary baseline technique to extract the target safety vector and test its geometric relationship to utility.

\paragraph{Assessing the Brittleness of Safety Alignment via Pruning and Low-Rank Modifications (Wei et al., 2024).}
\citet{wei2024brittleness} expanded on the geometry of safety by examining its broader representational footprint. They proved that safety mechanisms occupy a remarkably sparse low-rank subspace within the overall parameter space of a model. By applying Activation-aware Singular Value Decomposition (actSVD) to the internal activations, they successfully separated the safety-critical regions from the utility-critical regions. Astonishingly, they found that these safety regions comprised only about 2.5\% of the rank level. They demonstrated that pruning these specific low-rank regions compromised safety without triggering a collapse of the general utility of the model. We utilize this actSVD methodology as our second extraction and ablation technique. Comparing the sparse subspace identified by actSVD against the single vector extracted via DoM allows us to robustly map the boundaries between safety and utility representations.
